\scanpage{168}

\begin{center}Unités.

---\end{center}

\uline{Soit} $K$ \uline{un corps de nombres algébriques plongé dans} $\mathbf C$. \uline{Les conditions suivantes sont alors équivalentes} : a) \uline{pour toute unité} $\varepsilon$ \uline{de} $K$, \uline{il y a} $n\in\mathbf N$, $n>0$ \uline{tel que} $\varepsilon^n$ \uline{soit réel} ; b) \deleted{$\unsure{K\subset\mathbf R}$, \unsure{ou}} \uline{ou bien} $K\subset\mathbf R$ \uline{ou bien} $K$ \uline{se met sous la forme} $K_0(\alpha)$, \uline{avec} $K_0$ \uline{totalement réel}, $\alpha^2\in K_0$, $\alpha^2$ \uline{totalement négatif}.

Supposons $K\not\subset\mathbf R$ et soit $K'=K\cap\mathbf R$ ; $K$ est extension \deleted{quadratique} \inserted{de degré $p\geq 2$} de $K'$. Soient $r_1$ (resp. $r'_1$) le nombre de conjugués réels de $K$, $r_2$ (resp. $r'_2$) la moitié du nombre de conjugués imaginaires de $K'$ ; on a $[K:K']=2$, d'où $r_1+2r_2=p(r'_1+2r'_2)$. Si a) est satisfaite, le groupe des unités de $K'$ a même rang que celui de $K$, d'où \deleted{$\unsure{r_1+r_2}$} $r_1+r_2=r'_1+r'_2$ ; il en résulte que $0\geq r_1+2r'_2$, d'où $r_1=r'_2=0$ : $K'$ est totalement réel et $K$ totalement imaginaire. Réciproquement, si b) est satisfaite, on a $r_1=r'_2=0$, $r_2=r'_1$, $r_1+r_2=r'_1+r'_2$, et le groupe des unités de $K'$ a même rang que celui de $K$, d'où a).

\uline{Exemple}. La condition a) est satisfaite si $K=\mathbf Q(\unsure{\zeta})$, $\zeta^m=1$, car, $K/\mathbf Q$ étant abélienne, le corps des nombres réels de $K$ est totalement réel tandisque $K$ est totalement imaginaire.

\margin{167}
